\documentclass[10pt]{article}
\usepackage[T1]{fontenc}
\usepackage[utf8]{inputenc}
\usepackage[french]{babel}
\usepackage{amsmath,amssymb,mathrsfs,mathpazo}
\usepackage{geometry}
\geometry{a4paper,margin=2cm}
\usepackage{graphicx,array,tabularx,tabularray,colortbl,multirow,multicol,booktabs,enumitem,pifont,ulem,fancybox,diagbox,cancel}
\usepackage{pstricks,pstricks-add,pst-plot,pst-tree,pst-node,pst-text,pst-func,pst-math,pst-bspline,pst-3dplot,pst-eucl,pst-solides3d}
\let\pscalculate\undefined
\usepackage{pst-bezier}
\usepackage{tikz,pgfplots,tkz-tab}
\usetikzlibrary{arrows,arrows.meta,patterns,calc,positioning,trees,decorations.markings,shapes.geometric,shapes.misc,shapes.arrows,intersections,decorations.pathmorphing,decorations.pathreplacing,backgrounds,shadows,babel}
\pgfplotsset{compat=1.18}
\usepgfplotslibrary{fillbetween}
\usepackage[np]{numprint}
\usepackage{listings,makecell,pas-tableur,twoopt}
\let\barre\undefined
\usepackage[pstricks]{tablvar}
\definecolor{darkgreen}{RGB}{0,100,0}
\providecommand{\red}{\color{red}}
\providecommand{\blue}{\color{blue}}
\providecommand{\green}{\color{green}}
\providecommand{\black}{\color{black}}
\providecommand{\white}{\color{white}}
\providecommand{\gray}{\color{gray}}
\providecommand{\cyan}{\color{cyan}}
\providecommand{\magenta}{\color{magenta}}
\providecommand{\yellow}{\color{yellow}}
\providecommand{\violet}{\color{violet}}
\providecommand{\orange}{\color{orange}}
\providecommand{\euro}{\texteuro}
\providecommand{\EUR}{\texteuro}
\providecommand{\ieme}{\textsuperscript{e}}
\providecommand{\degres}{\ensuremath{^\circ}}
\providecommand{\degree}{\ensuremath{^\circ}}
\providecommand{\vv}[1]{\overrightarrow{#1}}
\providecommand{\Coord}[1]{#1}
\providecommand{\decosix}{}
\providecommand{\starredbullet}{\ensuremath{\bullet}}
\pagestyle{empty}
\setlength{\parindent}{0pt}
\usepackage{adjustbox}\DeclareUnicodeCharacter{2236}{:}\DeclareUnicodeCharacter{2190}{\ensuremath{\leftarrow}}
\providecommand{\R}{\mathbb{R}}
\providecommand{\N}{\mathbb{N}}
\providecommand{\Z}{\mathbb{Z}}
\providecommand{\Q}{\mathbb{Q}}
\providecommand{\D}{\mathbb{D}}
\providecommand{\C}{\mathbb{C}}
\providecommand{\e}{{\mathrm{e}}}
\providecommand{\vect}[1]{\overrightarrow{#1}}
\providecommand{\vectt}[1]{\overrightarrow{\mathrm{#1}}}
\providecommand{\vvt}[1]{\overrightarrow{\mathrm{#1}}}
\providecommand{\np}[2][]{{#2~#1}}
\providecommand{\vv}[1]{{\overrightarrow{#1}}}
\renewcommand{\d}{{\mathrm{d}}}
\providecommand{\degres}{{^\circ}}
\providecommand{\degre}{{^\circ}}
\providecommand{\blue}{\color{blue}}
\providecommand{\red}{\color{red}}
\providecommand{\enskip}{\quad}
\providecommand{\llbracket}{\lbrack\!\lbrack}
\providecommand{\rrbracket}{\rbrack\!\rbrack}
\providecommand{\euro}{{\text{€}}}
\providecommand{\no}{{\text{n°}}}
\providecommand{\interval}[2]{{[#1;#2]}}
\providecommand{\intervFO}[2]{{[#1;#2[}}
\providecommand{\dotfill}{\ldots}

\providecommand{\ds}{\displaystyle}
\providecommand{\textding}[1]{#1}
\providecommand{\Oij}{\(\mathrm{O};\vec{\imath},\vec{\jmath}\)}
\providecommand{\Oijk}{\(\mathrm{O};\vec{\imath},\vec{\jmath},\vec{k}\)}
\providecommand{\Ouv}{\(\mathrm{O};\vec{u},\vec{v}\)}
\providecommand{\intervFF}[2]{\ensuremath{[#1;#2]}}
\renewcommand{\euro}{\ifmmode\text{\texteuro}\else\texteuro\fi}
\DeclareUnicodeCharacter{00B0}{\ifmmode^{\circ}\else\textdegree\fi}
\DeclareUnicodeCharacter{20AC}{\ifmmode\text{\texteuro}\else\texteuro\fi}
% Énoncés : Éducation nationale. Transcription et corrigés : APMEP, Denis Vergès et contributeurs. Découpage et classement : CoopMaths, Sébastien Lozano, Éric Elter et contributeurs. Adaptation de la présentation : Maths Explorer.
% CC BY-SA : https://coopmaths.fr/www/about/
% https://coopmaths.fr/alea/static/bac/2025/tex/bac_2025_05_sujet2_ameriquenord_secours_3.tex
\begin{document}
\noindent\textit{APMEP / CoopMaths. CC BY-SA. Présentation adaptée par Maths Explorer.}\par\medskip
\medskip

Dans cet exercice, les réponses seront arrondies à $10^{-4}$ près.

\bigskip

Durant la saison hivernale, la circulation d'un virus a entraîné la contamination de $2\,\%$ de la population d'un pays. Dans ce pays, $90\,\%$ de la population a été vaccinée contre ce virus.

On constate que $62\,\%$ des personnes contaminées avaient été vaccinées.


\bigskip

On interroge au hasard une personne, et on note les évènements suivants :

\begin{description}
\item[ ] $C$ : \og la personne a été contaminée\fg{}
\item[ ] $V$ : \og la personne a été vaccinée\fg{}.
\end{description}

Les évènements contraires des évènements $C$ et $V$ sont notés respectivement $\overline{C}$ et $\overline{V}$.

\medskip

\begin{enumerate}
\item À partir de l'énoncé, donner, sans calcul, les probabilités $P(C)$, $P(V)$ et 
la probabilité conditionnelle $P_C(V)$.
\item
	\begin{enumerate}
		\item Calculer $P(C \cap V)$.
		\item En déduire $P\left(\overline{C} \cap V\right)$.
	\end{enumerate}
\item Recopier l'arbre des probabilités ci-dessous et le compléter.

\begin{center}
\begin{tikzpicture}[xscale=1,yscale=1,baseline={(R.base)}]

\tikzstyle{fleche}=[->,>=latex,thick]
\tikzstyle{noeud}=[fill=white,circle,inner sep=2pt]
\tikzstyle{feuille}=[fill=white,circle,inner sep=2pt]
\tikzstyle{etiquette}=[pos=0.6,fill=white, inner xsep=3pt, inner ysep=1.5pt]

\def\DistanceInterNiveaux{3}
\def\DistanceInterFeuilles{0.8}

\def\NiveauA{(0)*\DistanceInterNiveaux}
\def\NiveauB{(1)*\DistanceInterNiveaux}
\def\NiveauC{(2)*\DistanceInterNiveaux}
\def\InterFeuilles{(-1)*\DistanceInterFeuilles}

\node[noeud] (R) at ({\NiveauA},{(1.5)*\InterFeuilles}) {};
\node[noeud] (Ra) at ({\NiveauB},{(0.5)*\InterFeuilles}) {$C$};
\node[feuille] (Raa) at ({\NiveauC},{(0)*\InterFeuilles}) {$V$};
\node[feuille] (Rab) at ({\NiveauC},{(1)*\InterFeuilles}) {$\overline{V}$};
\node[noeud] (Rb) at ({\NiveauB},{(2.5)*\InterFeuilles}) {$\overline{C}$};
\node[feuille] (Rba) at ({\NiveauC},{(2)*\InterFeuilles}) {$V$};
\node[feuille] (Rbb) at ({\NiveauC},{(3)*\InterFeuilles}) {$\overline{V}$};

\draw[fleche] (R.east)--(Ra.west);
\draw[fleche] (Ra.east)--(Raa.west);
\draw[fleche] (Ra.east)--(Rab.west);
\draw[fleche] (R.east)--(Rb.west);
\draw[fleche] (Rb.east)--(Rba.west);
\draw[fleche] (Rb.east)--(Rbb.west);
\end{tikzpicture}
\end{center}

\item Calculer $P_V(C)$ et interpréter le résultat dans le contexte de l'exercice.
\item Déterminer si les affirmations suivantes sont vraies ou fausses en justifiant votre réponse.
	\begin{enumerate}
		\item \og Parmi les personnes non contaminées, il y a dix fois plus de personnes vaccinées que de personnes non vaccinées.\fg

		\item \og Plus de 98\,\% de la population vaccinée n'a pas été contaminée.\fg{}
	\end{enumerate}
\item On s'intéresse à un échantillon de $20$ personnes choisies au hasard dans la population.

La population du pays est assez importante pour qu'on puisse assimiler ce choix à des tirages successifs avec remise.

On note $X$ la variable aléatoire qui à chaque tirage associe le nombre de personnes contaminées.

\emph{On rappelle que, pour une personne choisie au hasard, la probabilité d'être contaminée est} $p = 0,02$.
	\begin{enumerate}
		\item Quelle est la loi suivie par la variable aléatoire $X$ ? Justifier et donner ses paramètres.
		\item Calculer, en rappelant la formule, la probabilité que 4 personnes exactement soient contaminées dans ce groupe de 20 personnes.
	\end{enumerate}
\end{enumerate}

\bigskip

\end{document}
